\section{Introduction}\label{sec:intro}

\subsection{The input and the two questions}

Throughout, the input theorem is \cref{src:thm:dgrh}: the Riemann hypothesis for $\zeta(s)$
and for every Dirichlet $L$-function, in the form exported by the Lean development \repo\ at
the commit named in \cref{sec:src}. We refer to it as Dirichlet GRH. The article does not
re-verify it. Every result below that depends on it says so.

Two questions organize the material. The first is what Dirichlet GRH settles outside its own
statement, either through published theorems that assume it or through new arguments. The
second is which parts of its proof can be reused. \Cref{src:sec:mechanism} describes that
proof as the reports read it: a comparison of two bounds for the same covolume, a spanning
statement, a twisting step that selects torsion points with controlled valuations, and a
final analytic step that turns a one-sided bound for a smoothed sum over primes into the
absence of zeros. Each of these ingredients gives rise to one or more sections.

\subsection{Status of the results}

Three kinds of statement occur, and the optional argument of each theorem says which one it
is.
\begin{itemize}[leftmargin=2em]
\item \emph{\stU.} The statement is proved without any hypothesis on zeros. All of
  Part~V and of \cref{sec:spn}, and most of Part~VI, are of this kind.
\item \emph{\stG.} The proof uses \cref{src:thm:dgrh} and nothing else that is unproved.
\item \emph{\stR.} The statement is an equivalence or an implication between a Riemann
  hypothesis and an explicit inequality, usually a one-sided bound for a weighted count of
  rational primes. The inequality is not proved. No result of this kind establishes a new
  Riemann hypothesis, Artin's primitive-root conjecture, or any other named conjecture.
\end{itemize}
Finite computations enter several proofs. Where a result rests on one, the statement says
so and names the delivered file that records it. File names are relative to the directory
\begin{center}
\code{Analysis/GeneralizedRiemannHypothesis/Research}
\end{center}
of the repository that holds this article, as it stands at commit \code{10802a6b36} for the
103 reports of the first two batches, at commit \code{f2f9fd7531} for the 17 of the third, at
commit \code{64f3896557} for the 11 of the fourth, at commit \code{fd9a9cb2bd} for the 11 of the
fifth, at commit \code{f6b5955e13} for the 42 of the sixth and at commit \code{7f951e1dfd} for the 10 of
the seventh and at commit \code{3c4d27a69c} for the 15 of the eighth; the packet names are unique. Nothing was recomputed for the
merge, and the mathematical claims are those of the reports.

\subsection{Guide to results}\label{intro:sec:guide}

The table lists topics and questions a reader may look for, with the section that treats each
and, where it helps, the subsections that hold the main statements.

\begingroup\small
\renewcommand{\arraystretch}{1.12}
\begin{longtable}{@{}>{\raggedright\arraybackslash}p{0.58\textwidth}>{\raggedright\arraybackslash}p{0.38\textwidth}@{}}
\toprule
Topic or question & Where\\\midrule
\endfirsthead
\toprule
Topic or question & Where\\\midrule
\endhead
\bottomrule
\endfoot
\multicolumn{2}{@{}l}{\emph{Part I: the input, the source mechanism and the analytic engine}}\\
The input theorem: what exactly is assumed &
  \cref{sec:src}: \labelcref{src:sec:input}\\
Architecture of the source proof; quantifiers of the twisting step; what a transfer must
supply & \cref{sec:src}: \labelcref{src:sec:mechanism,src:sec:quantifiers,src:sec:limits}\\
The determinant comparison; paying for missing dimensions &
  \cref{sec:src}: \labelcref{src:sec:comparison}; \cref{sec:spn}: \labelcref{spn:sec:missing}\\
The spanning lemma at sharp and at the source's constants; finite selection certificates; why
almost all of the target must be spanned &
  \cref{sec:spn}: \labelcref{spn:sec:spanning,spn:sec:source-improvement,spn:sec:certificates,src:sec:span}\\
Integral frames, integral Gram determinants, self-dual projective lattices &
  \cref{sec:spn}: \labelcref{spn:sec:frames,src:sec:gram,spn:sec:projective}\\
One-sided Mellin--Landau zero tests; ramp sums; removing prime powers &
  \cref{sec:ana}: \labelcref{ana:sec:basic,ana:sec:ramp,ana:sec:primeonly}\\
Finite-difference, signed-mean and joint detectors; sharp constants for dyadic filters &
  \cref{sec:ana}: \labelcref{ana:sec:fd,ana:sec:mean,ana:sec:dyadic}\\
Constant-free statistics for families of Dirichlet characters &
  \cref{sec:ana}: \labelcref{ana:sec:family}\\
What averaged bounds recover pointwise; sampling a ramp on sparse scales; sampling meshes &
  \cref{sec:ana}: \labelcref{ana:sec:mpt,nab:sec:sampling,ana:sec:mesh}\\
Auxiliary windows; tests on thin sets of twists and of scales &
  \cref{sec:ana}: \labelcref{ana:sec:window,ana:sec:sparse}\\
Li coefficients: detection schedules, second differences, the character from filtered Li
coefficients & \cref{sec:li}: \labelcref{li:sec:sampling,rec:sec:li,li:sec:curv}\\
Moment pencils, Ritz values, matched ordinates, Hermite and Hankel power-sum tests, cluster
counts & \cref{sec:li}: \labelcref{li:sec:ritz,li:sec:match,li:sec:hermite,ana:sec:hankel,li:sec:cluster}\\
Even central derivatives of quadratic Dirichlet $L$-functions (Sheth and Voo) &
  \cref{sec:li}: \labelcref{li:sec:central}\\
The Li heat semigroup and L\'evy processes from Cayley angles; the character from one kernel,
one path or one time slice &
  \cref{sec:lih}: \labelcref{li:sec:univ,rec:sec:liheat,li:sec:path,li:sec:offorigin,li:sec:slice}\\
A small-deviation law for sums over zeros; Euler's Gamma product &
  \cref{sec:lih}: \labelcref{ana:sec:smalldev,ana:sec:eulergamma}\\
Best underapproximations by distinct unit fractions & \cref{sec:lih}: \labelcref{lih:sec:egypt}\\[3pt]
\multicolumn{2}{@{}l}{\emph{Part II: nonabelian zero detection}}\\
The Riemann hypothesis for the Dedekind zeta function of a Galois field from one prime
inequality; Heilbronn domination & \cref{sec:nab}: \labelcref{nab:sec:heilbronn,nab:sec:sylow}\\
Symmetric groups, reflection detectors, normal towers, deficits and nonsplit primes, the zeros
missing from the abelian part &
  \cref{sec:nab}: \labelcref{nab:sec:sn,nab:sec:reflection,nab:sec:tower,nab:sec:deficit,nab:sec:missing}\\
Smaller zero-detecting fields for $S_n$ & \cref{sec:nab}: \labelcref{nab:sec:lowdegree}\\
Which Frobenius weights can detect zeros; why fixed congruence weights cannot &
  \cref{sec:nab}: \labelcref{nab:sec:extremal,nab:sec:periodic}\\
Detectors for $A_5$, $A_6$, $\PSL_2(7)$, $\PSL_2(q)$, $S_4$, six quintic fields and the real
$S_3$ field of discriminant $229$ & \cref{sec:grp}; for $S_n$, \cref{nab:sec:sn}\\
Kummer fields and Artin's primitive-root conjecture &
  \cref{sec:grp}: \labelcref{grp:sec:kummer}\\
The cubic field of discriminant $-23$: prime races, eta products, sign tests &
  \cref{sec:hec}: \labelcref{hec:ss:cubic,hec:ss:othercubic}\\
Dihedral conductors $47$ and $71$; eta products and Jensen polynomials; class-group and
Gaussian Hecke characters; modular forms and symmetric powers &
  \cref{sec:hec}: \labelcref{hec:ss:dihedral,hec:ss:jensen,hec:ss:classgroup,hec:ss:gaussian,hec:ss:sym}\\[3pt]
\multicolumn{2}{@{}l}{\emph{Part III: zeros and prime sums under Dirichlet GRH}}\\
Prime races: modulo four in a fixed window, weighted races in which central zeros help,
moving Gaussian windows & \cref{sec:zer}: \labelcref{zer:sec:window,hec:ss:central,zer:sec:gauss}\\
Euler products at the critical line; Suzuki's prime-power questions &
  \cref{sec:zer}: \labelcref{zer:sec:euler,zer:sec:suzuki}\\
Zhao's Mertens-error question for abelian fields & \cref{sec:zer}: \labelcref{zer:sec:zhao}\\
Covariance, oscillation and balancing of prime errors at several scales &
  \cref{sec:flu}: \labelcref{zer:sec:cov,zer:sec:balancing}\\
Peaks, sharp windows and excursions of normalized prime errors; the source remainder &
  \cref{sec:flu}: \labelcref{zer:sec:peaks}\\
Merging scales: stencil variances and multiplicities of zeros &
  \cref{sec:flu}: \labelcref{zer:sec:coalesce}\\
Shortest positive averages and filters that cancel prescribed zeros; logarithmic support &
  \cref{sec:fil}: \labelcref{fil:sec:reflect,fil:sec:logsupport,fil:sec:global}\\
Reading the central multiplicity from prime sums with one positive rule or a finite grid &
  \cref{sec:fil}: \labelcref{fil:sec:primes}, \cref{ana:thm:grid}\\
Zero-free intervals of trigonometric polynomials: the two-frequency classification, the gaps of
Kozma and Oravecz for progressions, squares and higher powers &
  \cref{sec:zfg}: \labelcref{zfg:sec:dual,fil:sec:two,zfg:sec:prog,zfg:sec:squares,zfg:sec:powers}\\
Steinerberger's support problem; reciprocal moments of zero heights &
  \cref{sec:zfg}: \labelcref{zfg:sec:recip}\\
Erd\'elyi's question on the Lorentz degree of polynomials with zeros in a thin ellipse &
  \cref{sec:zfg}: \labelcref{zfg:sec:lorentz}\\
Counting zeros from values to the right of the critical strip; certificates for central
nonvanishing & \cref{sec:rec}: \labelcref{rec:sec:count}\\
Zero multisets and multiplicities from prime sums, sparse samples, Taylor data and extremal
profiles & \cref{sec:rec}: \labelcref{rec:sec:taylor,rec:sec:prime,rec:sec:sparse,rec:sec:profile}\\
The modulus, the character and twists from the zero spectrum, spacings, blurred zeros, poles,
low zeros and balanced statistics & \cref{sec:chz}: \labelcref{rec:sec:spectrum,rec:sec:gaps,rec:sec:blur,rec:sec:poles,rec:sec:quadratic,rec:sec:lowzero,rec:sec:homometry}\\
Both signs of multiplicity differences, and ranks of multiplicity vectors, in short windows &
  \cref{sec:chz}: \labelcref{rec:sec:localwin,zer:sec:short}\\
Matching and binning the zeros of two products on a dyadic range &
  \cref{sec:chz}: \labelcref{chz:sec:matching}\\
Graph couplings of completed $L$-functions: real zeros, inverse problems, transport of zeros,
local Cauchy laws & \cref{sec:cpl}: \labelcref{zer:sec:graph,zer:sec:transport}\\
How often two products of Dirichlet $L$-functions share Euler factors &
  \cref{sec:eul}: \labelcref{rec:sec:qmatch,rec:sec:cmatch,rec:sec:omatch}\\[3pt]
\multicolumn{2}{@{}l}{\emph{Part IV: arithmetic consequences of Dirichlet GRH}}\\
Consequences of published conditional theorems: abelian fields, least primes in progressions;
primitive roots and Euclidean fields & \cref{sec:dir}: \labelcref{dir:sec:published,dir:sec:roots}\\
A prime plus a multiple of a prime cube; primes in geometric intervals; prime-generating
constants; prime-gap graphs &
  \cref{sec:dir}: \labelcref{dir:sec:cube,dir:sec:geometric,dir:sec:mills,dir:sec:graphs}\\
Squarefree products of small primes in residue classes; factorizations, disjoint products,
transition profiles, routing &
  \cref{sec:prd}: \labelcref{dir:sec:products,dir:sec:factorizations,dir:sec:disjoint,dir:sec:profiles,dir:sec:routing}\\
Sign patterns of characters at small primes; prime codes and designs &
  \cref{sec:prd}: \labelcref{dir:sec:signs,dir:sec:codes}\\
Class-group exponents of imaginary quadratic fields and orders: Euler's list, exponents up to
eight; representations $n=xy+yz+zx$ &
  \cref{sec:cls}: \labelcref{cls:sec:inputs,dir:sec:idoneal,dir:sec:xyz,cls:sec:six,cls:sec:seven}\\
Coefficient fields of CM newforms, reflected fields, isogeny bounds, heavenly abelian varieties &
  \cref{sec:cmf}: \labelcref{dir:sec:cm,dir:sec:reflected,dir:sec:isogeny,dir:sec:heavenly}\\
Sparse polynomial equations for singular moduli; growth of the least support &
  \cref{sec:smd}: \labelcref{smd:sec:classification,smd:sec:growth,smd:sec:uncond}\\
$p$-adic fewnomial root counts, also beyond $p$ terms &
  \cref{sec:smd}: \labelcref{smd:sec:disk,smd:sec:exact,smd:sec:beyond}\\[3pt]
\multicolumn{2}{@{}l}{\emph{Part V: combinatorial transfers}}\\
Sets without three-term geometric progressions in $\F_q[x]$ and in finite coefficient rings &
  \cref{sec:gpp}: \labelcref{gpp:sec:counterexample,gpp:sec:upper,gpi:sec:poly,gpp:sec:rings,gpp:sec:squarezero}\\
The same over the integers and the ideals of number fields (McNew's questions) & \cref{sec:gpi}\\
Additive constraints from the sparse-place construction: periodic sets, anticodes &
  \cref{sec:gpp}: \labelcref{gpp:sec:syndetic,gpp:sec:sparse,gpp:sec:lattice}\\
$p$-adic valuations of spanning-tree counts in graph towers & \cref{sec:gra}\\
Periodic quantum walks, oriented Cayley walks, signed square cycles, Ehrhart zeros &
  \cref{sec:grw}: \labelcref{gra:sec:grover,gra:sec:cayley,gra:sec:cycles,gra:sec:ehrhart}\\
Ehrhart theory of graph polytopes: reconstruction from multipartite symmetric edge polytopes,
complete graphs with a matching deleted, gamma vectors of suspensions, the Davis--Kohl conjecture
for stable set polytopes &
  \cref{sec:grw}: \labelcref{grw:sec:multipartite,grw:sec:matching,grw:sec:gamma,grw:sec:stable}\\
Laws of compressions of random symplectic and orthogonal matrices & \cref{sec:cmp}\\
Relation spaces of structured random constraints, products of random codes, Schur squares &
  \cref{sec:cmp}: \labelcref{spn:sec:relations}\\[3pt]
\multicolumn{2}{@{}l}{\emph{Part VI: local transfers and the arithmetic interface}}\\
Torsion specialization: the optimal order of the selected root of unity and the exact minimax
cost & \cref{sec:tors}: \labelcref{cost:sec:minimax,tors:sec:optimal}\\
Blocking finite sets, torsion shells, ramified bases, nonvanishing alone &
  \cref{sec:tors}: \labelcref{cost:sec:finite,cost:sec:shell,cost:sec:ramified,tors:sec:nonvanishing}\\
Separate thresholds, packets of fixed cardinality, positive tolerance &
  \cref{sec:tors}: \labelcref{cost:sec:packet,cost:sec:cardinality,tors:sec:tolerance}\\
Prescribed coefficient fields, simultaneous powers, matrices and Fitting ideals &
  \cref{sec:cost}: \labelcref{tors:sec:fixedfield,cost:sec:fields,cost:sec:abelian,cost:sec:radical,tors:sec:powers,tors:sec:matrices}\\
From torsion tuples to Dirichlet and anticyclotomic characters; conductor bounds &
  \cref{sec:cost}: \labelcref{tors:sec:characters,spn:sec:characters,spn:sec:conductor,spn:sec:benchmark}\\
Hypotheses of the selection theorems that cannot be dropped &
  \cref{sec:cost}: \labelcref{cost:sec:boundary,tors:sec:limits}\\
Lubin--Tate branches, pole bounds, jet recovery, Fleck congruences & \cref{sec:lub}\\
The Selmer interface, trace depth, odd cocycle modules, annihilator witnesses, K-theory and
Vandiver's conjecture &
  \cref{sec:sel}: \labelcref{sel:sec:interface,sel:sec:tracedepth,sel:sec:oddsha,spn:sec:annihilator,sel:sec:ktheory}\\[3pt]
\multicolumn{2}{@{}l}{\emph{Throughout}}\\
Which results use Dirichlet GRH, which are reductions, which use nothing &
  the mark in the heading of each statement\\
Open questions & the last subsection of every section (for \cref{sec:tors}:
  \cref{tors:sec:open} in \cref{sec:cost})\\
\end{longtable}
\endgroup

\subsection{Contents}

After this introduction the article has thirty sections in six parts.

\paragraph{Part I: the input, the source mechanism and the analytic engine.}
\Cref{sec:src} states the input and the architecture of its proof, proves the determinant
comparison, and records the quantifiers of the twisting step and what any transfer must supply.
\Cref{sec:spn} treats the spanning and completion step of the source and its integral
refinements: a sharp spanning interface, the source theorem at its original constants, exact
selection certificates, optimal integral frames, integral Gram determinants and self-dual
projective lattices; all of it is unconditional. \Cref{sec:ana} contains the analytic step
in the generality needed later. Its main statement is the one-sided test
\cref{ana:thm:onesided}; \cref{ana:thm:primeonly} removes prime powers, and
\cref{ana:thm:window} gives the auxiliary averaging window. The section also treats
detectors that need only a bound in mean, sharp constants for dyadic filters, constant-free
statistics for families of characters, the question of what an averaged bound recovers
pointwise, the classification of sampling meshes, and tests on thin sets of scales and of
twists. \Cref{sec:li} concerns Li coefficients: on which schedules they detect an off-line
zero, how filtered Li coefficients determine the character, how the moment pencil they define
recovers individual zeros and certifies their simplicity, and what near-positivity of Hermite
and Hankel power-sum matrices says about the location of zeros. \Cref{sec:lih} studies the
infinitely divisible laws defined by the Cayley angles of the zeros and the Li heat kernel,
which determine the character from one kernel, one path or one time slice, and proves a
small-deviation law for sums over zeros and a finite certificate for Euler's Gamma product.

\paragraph{Part II: nonabelian zero detection.}
\Cref{sec:nab} reduces the Riemann hypothesis for the Dedekind zeta function of an arbitrary
Galois extension of $\Q$ to one inequality over rational primes, given Dirichlet GRH
(\cref{nab:thm:sylow}). The main tool is a domination of zero multiplicities by Heilbronn
characters, \cref{nab:thm:heilbronn}, which avoids Artin's holomorphy conjecture. The section
also gives detectors for symmetric groups and reflection detectors, compresses along normal
towers, determines which Frobenius weights can serve as detectors and proves that weights
depending only on a fixed congruence cannot (\cref{nab:thm:periodic}). \Cref{sec:grp} makes
the reduction explicit for particular groups and fields, among them $A_5$, $A_6$, $\PSL_2(q)$,
$S_4$, six quintic fields, Kummer fields and the real $S_3$ field of discriminant $229$, with
optimal detectors in several cases. \Cref{sec:hec} treats $L$-functions of degree two: the
cubic field of discriminant $-23$, dihedral conductors $47$ and $71$, eta products and Jensen
polynomials, class-group and Gaussian Hecke characters, and symmetric powers.

\paragraph{Part III: zeros and prime sums under Dirichlet GRH.}
\Cref{sec:zer} collects sign results: which of two prime sums leads, when an Euler product
converges at the critical line, and the exact upper limits of Mertens-type errors, including
the classification of abelian fields in Zhao's question (\cref{zer:thm:zhao}).
\Cref{sec:flu} proves the positive definiteness and rigidity of the joint fluctuations of
smoothed prime errors at several scales, without any independence hypothesis on the zeros,
and studies simultaneous balancing, extreme values and merged scales. \Cref{sec:fil}
determines the shortest positive averages, and the least logarithmic support of positive
filters, that cancel prescribed zeros or every zero up to a given height, and reads the
central multiplicity from prime sums. \Cref{sec:zfg} treats zero-free intervals and positivity
for trigonometric and algebraic polynomials: the two-frequency classification, the gaps of Kozma
and Oravecz for progressions and squares, Steinerberger's support problem and Erd\'elyi's question
on the Lorentz degree; it is unconditional apart from applications to zeros of Dirichlet
$L$-functions. \Cref{sec:rec} asks what prime data, samples to the
right of the critical strip and local Taylor data determine about the zeros: it counts zeros
from values to the right of the strip and recovers zero multisets and multiplicities from
filtered, sparse and extremal prime sums. \Cref{sec:chz} treats the opposite direction: it
recovers the modulus, the character, twists, omitted primes, products and abelian fields from
the zero spectrum, zero spacings, zeros in short windows, blurred zero measures, poles of prime
series and finitely many low zeros. \Cref{sec:cpl} couples completed Dirichlet $L$-functions
along a finite graph, proves that admissible couplings have only real zeros, solves the
inverse problem, and describes the transport of zeros with local Cauchy laws for logarithmic
derivatives. \Cref{sec:eul} determines how often two products of Dirichlet $L$-functions can
have the same Euler factor at a prime.

\paragraph{Part IV: arithmetic consequences of Dirichlet GRH.}
\Cref{sec:dir} collects consequences obtained through published conditional theorems and
direct ones for primitive roots, Euclidean fields, primes in geometric intervals,
prime-generating constants and prime-gap graphs. \Cref{sec:prd} studies squarefree products
of small primes in residue classes, and small primes as samples of characters: sign patterns,
codes and designs. \Cref{sec:cls} proves that the known lists of imaginary quadratic fields of
class-group exponent six and seven are complete (\cref{cls:thm:six,cls:thm:seven}), and with
Euler's list holds every class-group exponent census of the article. \Cref{sec:cmf} bounds
coefficient fields of CM newforms, reflected fields, isogenies of two families of sextic
fields and heavenly abelian varieties. \Cref{sec:smd} studies the least number of terms of a
polynomial that vanishes at a singular modulus; its classification of equations with at most
five terms (\cref{smd:thm:five}) is unconditional, and Dirichlet GRH enters only in a growth
estimate (\cref{smd:thm:growth}).

\paragraph{Part V: combinatorial transfers.}
The source proof uses additive constraints on exponents to exclude multiplicative patterns.
\Cref{sec:gpp,sec:gpi} apply the same idea to sets without three-term geometric progressions:
over polynomial rings and finite coefficient rings, where a published density conjecture is
disproved (\cref{gpp:thm:counterexample}), and over the integers and the ideals of a number
field, where several questions of McNew are answered
(\cref{gpi:thm:integer-gap,gpi:thm:beta-alpha}). \Cref{sec:gra} concerns the valuation of
the number of spanning trees in towers of graphs. \Cref{sec:grw} settles four spectral
problems on finite graphs, namely periodic Grover walks, least periods of oriented Cayley walks,
optimal signings of square cycles and the zeros of a family of Ehrhart polynomials, and carries
the Ehrhart theory of graph polytopes further: complete multipartite symmetric edge polytopes and
the reconstruction of the graph, complete graphs with a matching deleted, face realizations of
gamma vectors of graph suspensions, and the Davis--Kohl conjecture on stable set polytopes of
complements of odd cycles.
\Cref{sec:cmp} determines the laws of compressions of random symplectic and orthogonal
matrices over finite fields and local rings, which control the compressions used by the
spanning step.

\paragraph{Part VI: local transfers and the arithmetic interface.}
The twisting step of the source proof chooses a root of unity at which finitely many $p$-adic
power series keep a prescribed valuation. \Cref{sec:tors} proves the selection theorem with
the optimal bound on the order of the root (\cref{tors:thm:selection}) together with the
matching lower bounds, among them the exact cost of defeating every such choice
(\cref{cost:thm:minimax}). \Cref{sec:cost} determines the dependence on the coefficient field,
treats powers, matrices and characters, turns torsion data into global twists with explicit
conductor bounds, and shows which hypotheses cannot be dropped. \Cref{sec:lub} proves sharp
precision statements for branches of Lubin--Tate formal groups (\cref{lub:thm:pole}).
\Cref{sec:sel} states what the Selmer-group argument of the source proof requires of a Galois
representation, proves the resulting obstructions (\cref{sel:thm:smith}), compares the
modules that occur with standard Galois cohomology, and determines the exact greatest common
divisor of the cyclotomic annihilator witnesses (\cref{spn:thm:gcd}).

\subsection{How the reports were merged}

The reports were written in sequence, and later ones often repeat, sharpen, correct or
answer earlier ones. Each result appears here once. When several reports prove the same
statement, one proof is kept. When a result follows from another result of the collection
whose hypotheses are no stronger, it is omitted, or kept as a corollary or an example if it
carries an explicit constant or a worked instance. When a later report answers a question
of an earlier one, the answer is stated and the question is dropped; each section ends with
the questions that remain open. Descriptions of the source proof, which most reports
contain, are replaced by \cref{sec:src}. A line under each subsection heading names the
reports it is drawn from, and the tables in the directory \code{provenance} beside this
article list, for every section of every report, where it went and, for each omitted
result, the result that supersedes it.
